Pile-up (i.e. event overlap) contamintaion was evaluated in \href{https://phenix-intra.sdcc.bnl.gov/phenix/WWW/p/draft/mlario/phi_in_HeAu_pAu_pAl_v6.pdf}{AN1399} with the use of BBCLL1 narrow to Clock trigger rate. 
The value was found to be 1\% in MB, with no account to centrality. 
The determination procedure of the collision in the central region is the most prone to overlap contamination. 
Therefore pile-up rejection cut was introduced using $frac$ variable from the module DoubleInteractionUtils described in \href{https://phenix-intra.sdcc.bnl.gov/phenix/WWW/p/draft/dcm07e/DoubleInteractionAN/small-system-double.pdf}{AN1304}.

The frac distribution dependence on centrality is shown as an occupancy histogram on Fig. \ref{fig:frac_cut}. 
High occupancy regions are peaks on frac distribution on each centrality slice. 
These peaks mostly come from clean (or cleanly reconstructed) events. 
Overlapping events span almost the whole frac region. 
In \href{https://phenix-intra.sdcc.bnl.gov/phenix/WWW/p/draft/dcm07e/DoubleInteractionAN/small-system-double.pdf}{AN1304} it was shown that removing events with frac values outside the peaks removes most of overlap contamination. 
Therefore the centrality dependent cut with a function of a form $\text{frac}_{\text{threshold}} = p_0 - p_1 e^{p_2 \cdot \text{centrality}^2}$ was introduced with parameters $p_0 = 80$, $p_1 = 0.5$, $p_2 = 5.5$. 
This function is shown as a red line of Fig. \ref{fig:frac_cut}. Events to the left of $\text{frac}_{\text{threshold}}$ were removed from the analysis.
The choice of $\text{frac}_{\text{threshold}}$ function was due to its shape closely following selected frac threshold values for all centrality slices.

\begin{figure}[h!]
	\centering
	\includegraphics[width=0.5\textwidth]{FracFit/Run14HeAu200.pdf}
   \caption[$\text{frac}$ vs centrality distribution and $\text{frac}_{\text{threshold}}$ cut as a function of centrality]{$\text{frac}$ vs centrality distribution and $\text{frac}_{\text{threshold}}$ cut as a function of centrality (red line).}
	\label{fig:frac_cut}
\end{figure}
